\chapter{Introduction}
\label{chp:HCIntro}
\fancyhead[C]{Chapter \ref{chp:HCIntro}}

\section{Steady State Heating}
\label{sec:HCSSHeatLoss}

There are several steps to calculate the required heating for a room and the size of the central plant that is required, the general procedure is as follows:

\begin{itemize}
\item Calculate Conduction (Fabric) Heat Loss,
\item Calculate Infiltration Heat Loss,
\item Calculate Ventilation System Heat Loss (If Applicable),
\item Allow for Local Effects,
\item Determine Intermittent Correction Factor,
\item Calculate Peak Heating Load,
\item Size/Select Heat Emitters,
\item Calculate Distribution System Heat Loss,
\item Determine Infiltration Diversity Ratio,
\item Calculate Central Plant Size,
\item Select Central Plant, including any redundancy.
\end{itemize}

\subsection{Conduction (Fabric) Heat Loss}

\begin{equation}\label{eq:HeatUValue}
Q = U A
\end{equation}

Define U-value here...

\begin{equation}\label{eq:UValueDef}
U = \frac{1}{(\sum R + R_i + R_o)}
\end{equation}

Where R is the thermal resistance of the elements, R\textsubscript{i} and R\textsubscript{o} are the internal and external thermal resistances respectively.

{\href{https://www.designingbuildings.co.uk/wiki/U-values}{\underline{Designing Buildings Wiki on U-values}}

\subsection{Infiltration Heat Loss}

e.g. A air pressure test result of \SI{5}{\cubic\metre\per\square\metre\per\hour} at \SI{50}{\pascal} would result in a air infiltration rate of 0.111 ACH.

\subsection{Ventilation Heat Loss}



\subsection{Intermittent Correction Factor}



\subsection{Peak Heating Load}



\subsection{Distribution Losses}



\subsection{Infiltration Diversity Ratio}



\subsection{Central Plant Size}



\section{Annual Estimation of Heating Loads}